\section{Source estimates for the exterior curvature}
\label{gt-source-curvature}

Use \(t=\lambda\) for physical heat time. Fix a real center \(a\) and a
counted pair as in \eqref{eq:cluster-coordinates}--\eqref{eq:pair-discriminant}.
On each slice, let \(F\) be the field, \(Q(z)=(z-c)^2-D\) its pair
polynomial, \(A=F/Q\) the zero-free factor on the counted disk, and
\(B=A'/A\). The drift \(K\) is the divided difference of \(B\) at the
two roots, with \(B'(c)\) at a double. Multiplying the field by a
nonzero scalar on a slice leaves these quantities unchanged. The
physical drift enters \eqref{eq:gt-gap-law}, with the signed survival
budget given by \eqref{eq:gt-backward-gap}.

\subsection{Laurent coefficients of the complete field}

\begin{proposition}[Positive Laurent modes of the deflated logarithm]
\label{gt-laurent-modes}
Suppose that \(|z-a|<R\) contains the pair, the boundary is zero-free
with a collar, and the offsets \(u_i=z_i-a\) satisfy \(|u_i|\le r<R\).
Let \(L(w)\) be any collar logarithm of \(F(a+w)/w^2\), with a
constant centering allowed. Define
\[
 g_n=\frac{1}{2\pi\ii}\oint_{|w|=R}L(w)w^{-n-1}\,\dd w,\qquad
 h_j(u_1,u_2)=\sum_{\ell=0}^j u_1^\ell u_2^{j-\ell}.
\]
Then
\begin{equation}
 K=\sum_{n\ge2} n g_n h_{n-2}(u_1,u_2).
 \label{gt-laurent-k}
\end{equation}
The series converges absolutely. For a field real under conjugation,
\(g_n\) for \(n\ge1\) and the symmetric \(h_j\) are real.
\end{proposition}

\begin{proof}
The deflated factor is holomorphic and zero-free in a slightly larger
disk. Thus \(g(w)=\log A(a+w)\) exists there. For \(|w|>r\), expanding
both logarithms of \(1-u_i/w\) yields
\[
 L(w)=g(w)+C-\sum_{m\ge1}\frac{u_1^m+u_2^m}{m w^m}.
\]
Here \(C\) includes the branch difference and the centering. Its value
cannot affect the positive Laurent coefficients, which are the Taylor
coefficients of \(g\). Conjugation of the pair and field makes \(g'\)
real under conjugation, giving the asserted reality. Multiplying the
two convergent geometric series gives
\[
 \frac{1}{Q(a+w)}=\sum_{j\ge0} h_j(u_1,u_2)w^{-j-2}.
\]
In the contour formula \(K=(2\pi\ii)^{-1}\oint g'(w)/Q(a+w)\,\dd w\),
the \(w^{-1}\) coefficient pairs \(g'\)'s term \(n g_n w^{n-1}\)
with \(j=n-2\). Both series converge absolutely and uniformly on
the circle. At a double \(h_j=(j+1)u_1^j\), so the same formula
requires no division by a vanishing root separation.
\end{proof}

If \(|L|\le M_L\) on the circle, Cauchy's estimate gives
\(|g_n|\le M_L/R^n\), and \(|h_j|\le (j+1)r^j\).
The upper allowance obtained by retaining \(g_2\) exactly and
paying the rest solely by these two bounds is
\begin{equation}
 K\le U_2=2g_2+
 \frac{M_L}{R^2}\left(\frac{2}{(1-r/R)^3}-2\right).
 \label{gt-laurent-allowance}
\end{equation}
The parenthesis follows by differentiating the geometric series
twice, then removing the \(n=2\) term. With the fixed-tube choices
\(r=3/4,\ R=21/20,\ M_L=1737/1000\), one has
\(r/R=5/7\), \(M_L/R^2=386/245\), and the parenthesis is \(335/4\).
Consequently
\[
 U_2\ge -\frac{772}{245}+\frac{12931}{98}
       =\frac{63111}{490}>0 .
\]
Even exact knowledge of \(g_2\) cannot make this allowance certify
\(K\le-59/100\) at a single time with those parameters. The bound
is on the displayed allowance. To use the exact \(g_2\) for this
target, one needs
\(\sum_{n\ge3} n g_n h_{n-2}\le-59/100-2g_2\).
That signed tail may be smaller than the allowance through cancellation.

The unsigned local data can also coexist with either sign of the actual
drift in the class of analytic heat fields. Put \(s=t-T\),
\(p=1+4\kappa s\), \(b=\kappa/p\), and \(D=2sp\), and consider
\[
 F_\kappa(t,a+w)=P(w^2-D),\qquad P=p^{-5/2}\e^{bw^2}.
\]
On \(p>0\), the elementary derivatives are
\[
  P_t/P=-10b-4b^2w^2,\quad P_w/P=2bw,\quad
  P_{ww}/P=2b+4b^2w^2,\qquad D'=2+8bD .
\]
It follows that \((\partial_t F_\kappa+\partial_w^2F_\kappa)/P
=-8b(w^2-D)-D'+8bw^2+2=0\). The deflated logarithmic derivative
is \(2bw\), so \(K=2\kappa/p\). On \(|s|\le1/20\) and
\(0<|\kappa|\le1/20\), \(99/100\le p\le101/100\) and
\(|D|\le101/1000\). Both roots lie inside \(7/10\).
On \(7/10<|w|<11/10\) a common bound follows from the branch
\[
 L=-\tfrac52\log p+bw^2+\log(1-D/w^2).
\]
Indeed \(|\log p|\le1/99\), \(|bw^2|\le121/1980\), and
\(|D/w^2|<1/4\), so the logarithm series bounds \(|L|\) by
\(277/660<1/2\). The count, collar and unsigned bounds therefore
allow both signs. For \(t_-=T-1/20\) and \(t_+=T+1/20\),
\(\exp(4\int_t^{t_+}K)=[p(t_+)/p(t)]^2\). The exact signed margin
is this weight at \(t_-\) times \(D(t_-)\), namely
\(-p(t_+)^2/(10p(t_-))\). It is zero if the starting time is \(T\).
For \(\kappa\ne0\) these examples have order two, from the quadratic
exponential, and the shifted field need not be globally even in \(z\).

\subsection{Strip, order and the exterior sum}

For the canonical \(H_t\), reflection of \(\xi\) makes the time-zero
formula in Section~\ref{sec:source} equivalent to
\(H_0(z)=\xi((1+\ii z)/2)/8\). The classical strip
\(0\le\Re s\le1\) gives \(|\Im z|\le1\) for its zeros in these
variables \cite[\S25.10(i)]{dlmf}. De Bruijn's strip theorem for a
real even integrable Fourier source of decay
\(O(\e^{-|u|^\gamma})\), \(\gamma>2\), replaces the squared
half-width \(y_0^2\) by \(\max(y_0^2-2t,0)\) after multiplication
of the source by \(\e^{tu^2}\). The heat parameter here is the one
in \eqref{eq:field}. The source envelope \eqref{eq:tail} pays the
decay hypothesis. Thus on \(I=[0,1/10]\) every canonical zero
\(\varrho\) satisfies
\begin{equation}
 |\Im\varrho|\le\sqrt{1-2t}\le1 .
 \label{gt-canonical-strip}
\end{equation}
This is the unconditional global strip, with nonreal and repeated
zeros allowed. See Newman's Theorem~1 and the same-normalization
formulation in Polymath, \S1 and Theorem~3.2 \cite{newman,polymath}.

The same source estimate pays the growth needed for the zero product.
For \(|z|\le R\) on \(I\), the absolute integrand on \(u\ge0\) in
\eqref{eq:field} is bounded, up to a constant, by
\(\exp(9u+u^2/10+Ru-\pi\e^{4u})\). Absorb \(9u+u^2/10\) into a
constant plus \((\pi/2)\e^{4u}\), retain the integrable factor
\(\exp(-(\pi/4)\e^{4u})\), and maximize the remaining exponent.
For \(R\ge0\),
\[
 \sup_{u\ge0}\bigl(Ru-(\pi/4)\e^{4u}\bigr)
 \le \tfrac R4\log_+(R/\pi),
\]
where the right side is zero at \(R=0\). Hence
\(\log\max_{|z|\le R}|H_t(z)|=O(R\log(2+R))\) as \(R\to\infty\),
uniformly on \(I\). Positivity gives \(H_t(0)\ge H_0(0)>0\).
Jensen's formula on radius \(2R\) yields a zero count
\(O(R\log(2+R))\) inside radius \(R\), since every such zero
contributes at least \(\log 2\). Summing inverse squares over
dyadic annuli gives the uniform tail
\(O(\log(2+R)/R)\), with no evaluated constant needed for convergence.

\begin{proposition}[Exterior divisor formula]
\label{gt-exterior-proposition}
For a real entire field of order at most one on a counted pair disk,
\begin{equation}
 -K=\sum_{\varrho\ {\rm exterior}}
       \frac{1}{(z_1-\varrho)(z_2-\varrho)}
 \label{gt-exterior-k}
\end{equation}
with multiplicities. The sum converges absolutely and applies also at
a double by setting \(z_1=z_2\).
\end{proposition}
\begin{proof}
For order at most one, the bound \(\log M(R)=O(R^{1+\epsilon})\)
for \(0<\epsilon<1\) similarly gives inverse-square summability,
after removing any zero at the origin for Jensen's formula.
Hadamard factorization uses factors
\((1-z/\varrho)\e^{z/\varrho}\) for nonzero zeros and has a
prefactor \(z^m\e^{\alpha+\beta z}\).
In the canonical case the equivalent opposite-pair product is the
one used in \cite[proof of Proposition~3.1(i)]{polymath}. Logarithmic
differentiation gives the terms \(1/(z-\varrho)+1/\varrho\), a
constant \(\beta\), and the origin term \(m/z\). Subtracting \(Q'/Q\)
removes precisely the two counted occurrences. The remaining
constants cancel in the divided difference. Each remaining simple-pole
term contributes its negative reciprocal product, including \(m\)
copies of the term for an exterior zero at the origin. The finite
poles stay outside the counted disk, and the reciprocal-product tail
is \(O(|\varrho|^{-2})\) on compact bounded pair coordinates.
Termwise divided differences and the confluent derivative are also
justified by taking the contour integral before passing to the limit.
\end{proof}

\subsection{Pairing conjugate exterior zeros}

\begin{proposition}[Grouped exterior lower mass]
\label{gt-grouped-exterior}
On a slice satisfying Proposition~\ref{gt-exterior-proposition}, suppose
an auxiliary disk \(|z-a|<R_0\) counts only the target pair and has
zero-free boundary. Assume \(R_0>\delta\), \(|c-a|\le C,\ |D|\le d\)
and that
\(|\Im\varrho|\le\delta\) for the exterior zeros. If
\[
 X=\sqrt{R_0^2-\delta^2}-C>0,\qquad X^2>\delta^2+d ,
\]
then all exterior conjugate groups contribute positively to \(-K\).
For each selected disjoint conjugation-invariant multiset of \(n_j\)
exterior zeros with
\(L_j\le|\Re\varrho-c|\le U_j\) and \(L_j^2>\delta^2+d\), its
contribution is at least \(n_j m(L_j,U_j)\), where
\begin{equation}
 m(L,U)=\frac{L^2-\delta^2-d}
                  {(U^2+\delta^2+d)^2+4U^2\delta^2}.
 \label{gt-group-mass}
\end{equation}
\end{proposition}
\begin{proof}
Write \(\varrho=c+x+\ii y\). The product in \eqref{gt-exterior-k}
is \((x+\ii y)^2-D\). For a nonreal conjugate pair its reciprocal
and the conjugate reciprocal sum to \(2q\), where
\begin{equation}
 q(x,y,D)=\frac{x^2-y^2-D}{(x^2-y^2-D)^2+4x^2y^2}.
 \label{gt-conjugate-weight}
\end{equation}
A real root contributes \(q(x,0,D)\). Equal multiplicities at
conjugate zeros are included. Every exterior root satisfies
\(|\varrho-a|\ge R_0\), so
\(|x|\ge\sqrt{R_0^2-\delta^2}-C=X\).
The numerator in \eqref{gt-conjugate-weight} is positive.
Within a selected multiset it is at least \(L_j^2-\delta^2-d\),
and its denominator is at most
\((U_j^2+\delta^2+d)^2+4U_j^2\delta^2\).
Each occurrence contributes its share, and the absolutely convergent
sum permits omission of all unselected conjugate groups in this
lower bound for \(-K\).
\end{proof}

When \(F\) is even and the reflected pair is exterior, retain
\(-z_1,-z_2\) as one multiset. Provided \(c^2-D>0\), its contribution
is exactly
\[
 \frac{1}{(2z_1)(z_1+z_2)}+\frac{1}{(z_1+z_2)(2z_2)}
   =\frac{1}{2z_1z_2}=\frac{1}{2(c^2-D)} .
\]
This includes both occurrences at a double. For a positive center
\(a>C\) with \((a-C)^2>d\), the lower bound for this contribution
is \(\kappa_0=[2((a+C)^2+d)]^{-1}\).
Adding other disjoint selected multisets from
Proposition~\ref{gt-grouped-exterior} gives
\begin{equation}
 K\le-\kappa_0-\sum_j n_j m(L_j,U_j).
 \label{gt-mirror-allowance}
\end{equation}
A uniform bound requires the counts and ranges on every slice, with
each occurrence assigned once. The lower mass supplies an upper
bound on \(K\). In the all-real exterior class an upper endpoint
inverse-square sum has the opposite direction. Indeed, for real
separated \(z_1,z_2\) and every real exterior \(\varrho\) outside the
segment, \(p=z_1-\varrho\) and \(q=z_2-\varrho\) give
\(p^{-2}+q^{-2}-2/(pq)=(p-q)^2/(p^2q^2)\ge0\). Thus an upper
bound \(G_*\) for the sum of \(p^{-2}+q^{-2}\) implies
\(K\ge -G_*/2\).

\begin{corollary}[Exterior mass required on an isolated tube]
\label{gt-fixed-strip-tube}
Assume the hypotheses of Proposition~\ref{gt-grouped-exterior} and
the bound \eqref{gt-mirror-allowance} hold throughout a time interval
\([u,v]\) of length \(s\). To reach \(K\le-\kappa\) by this estimate,
the other selected groups must supply
\(\sum_j n_jm(L_j,U_j)\ge\kappa-\kappa_0\) on each slice.
If \(D(v)<2s-4\kappa_0s^2\), the signed cost allowance from the
reflected-pair bound alone exceeds the incoming gap.
\end{corollary}
\begin{proof}
The required additional mass follows directly from
\eqref{gt-mirror-allowance}. For the constant drift bound
\(K\le-\kappa_0\), the signed cost allowance in
\eqref{eq:gt-backward-gap} satisfies
\[
 2\int_u^v\e^{-4\kappa_0(v-t)}\,\dd t
   \ge 2s-4\kappa_0s^2,
\]
by \(\e^{-x}\ge1-x\). Other exterior mass can reduce the actual cost.
\end{proof}

The auxiliary disk and target coordinate bounds require a separate
full-slab enclosure. Appendix~\ref{sec:gt-source-certification}
records a fixed canonical instance. The geometric premises and
the signed magnitude requirement can be evaluated independently
of the direct early endpoint.

\subsection{Sharp rectangular weights and necessary counts}

\begin{proposition}[Sharp rectangular lower weight]
\label{gt-sharp-exterior}
Let \(\delta,d\ge0\), \(0<L\le U\), and \(L^2>\delta^2+d\).
On the rectangle \(L\le |x|\le U,\ |y|\le\delta,\ |D|\le d\),
the exact minimum of \(q\) in \eqref{gt-conjugate-weight} is
\begin{equation}
 \mu(L,U;\delta,d)
 =\min_{\substack{r\in\{L,U\}\\ e\in\{-d,d\}}}
 \frac{r^2-\delta^2-e}{(r^2-\delta^2-e)^2+4r^2\delta^2}.
 \label{gt-sharp-weight}
\end{equation}
In \eqref{gt-mirror-allowance} one may replace each \(m(L_j,U_j)\)
by \(\mu(L_j,U_j;\delta,d)\).
\end{proposition}

\begin{proof}
Write \(s=x^2,\ u=y^2,\ v=s-u-D>0\), so \(q=v/(v^2+4su)\).
At fixed \(s,D\), its \(u\)-derivative has numerator
\(v^2-4s(s-D)\) and positive denominator \((v^2+4su)^2\).
The numerator is at most \(-(s-D)(3s+D)<0\), since
\(v\le s-D\) and \(s>d\ge |D|\). The minimum uses \(u=\delta^2\).
For fixed \(s,u\), the derivative with respect to \(v\) has numerator
\(4su-v^2\). On a positive interval this can only change sign from
positive to negative, so the minimum over \(D\) is at an endpoint.
Finally, with \(u,D\) fixed, the \(s\)-derivative has numerator
\(4u(u+D)-v^2\), which is strictly decreasing in \(s\).
This leaves \(s=L^2,U^2\). Collapsed intervals simply repeat
corners. The estimate \eqref{gt-group-mass} bounds \(q\) below on the
same rectangle, so \(\mu\ge m\). Summing over the selected occurrences
in Proposition~\ref{gt-grouped-exterior} gives the claim.
\end{proof}

\begin{corollary}[Necessary counts for the width-one allowances]
\label{gt-rectangle-count}
Under Corollary~\ref{gt-fixed-strip-tube}, take \(\delta=1,\ d=1/5\)
and a reflected-pair lower allowance \(0<\kappa_0<3/10^9\).
A lower bound of \(59/100\) for \(-K\) using \(\kappa_0\) and the weights
\eqref{gt-group-mass} requires at least ten additional selected
occurrences. Using the sharp rectangular weights \eqref{gt-sharp-weight}
requires at least five additional selected occurrences.
\end{corollary}

\begin{proof}
Put \(h=1+d=6/5\). At fixed \(L\), the denominator of
\(m(L,U)\) increases with \(U\), so the largest possible weight has
\(L=U=\sqrt{s}\), \(s>h\). With \(r=s-h\), the weight becomes
\[
 \frac{r}{r^2+4(h+1)r+4h(h+1)}.
\]
Its derivative has the sign of \(4h(h+1)-r^2\), giving the maximum
\(M_{\rm sep}=1/[4(h+1+\sqrt{h(h+1)})]=5/(44+4\sqrt{66})\).
For the sharp weight, the admissible values \(y=1,\ D=d\) give an
upper cap by maximizing \(r/(r^2+4r+4h)\) over \(r>0\).
The derivative has the sign of \(4h-r^2\). The maximum is
\(M_{\rm rect}=1/[4(1+\sqrt h)]\), at \(s_*=h+2\sqrt h\).
To see that this cap is attained within the rectangular information
class, take \(L=U=\sqrt{s_*}\). At \(D=-d\) the numerator variable
in \(q\) rises from \(v=2\sqrt h\) to \(v+2d\). The difference
of the corresponding weights has the sign of
\(4s_*-v(v+2d)=4\sqrt h(2-d)>0\).
Equation~\eqref{gt-sharp-weight} therefore gives \(M_{\rm rect}\)
for this rectangle.

Since \((2031/250)^2<66\), \(M_{\rm sep}<625/9562\).
The exact inequality
\[
 \frac{59}{100}-\frac{3}{10^9}-\frac{5625}{9562}
 =\frac{8289985657}{4781000000000}>0
\]
and \(\kappa_0<3/10^9\) rule out the separated estimate with nine
or fewer additional occurrences. Since \(4M_{\rm rect}<1/2\), the
sharp rectangular estimate with four or fewer also falls short.
\end{proof}

These thresholds constrain the two displayed lower estimates with
width one and \(d=1/5\). A uniform canonical estimate using that
information must pass at \(t=0\), even when the strip narrows later.
The cap \(5M_{\rm rect}>25/42>59/100\) allows each selected rectangle
to use the same optimizing horizontal coordinate and leaves
sufficiency for five canonical occurrences open. An actual selection
must account for disjoint groups, conjugate occurrences, a shared
\(D\), and full-time locations. Omitted groups can add positive mass
and data specifying reality can improve the lower weight. An exact
multiplicity count of one in a conjugation-invariant exterior region
forces a real simple zero. If its horizontal distance is at most
\(U\), it then supplies \(1/(U^2+d)\) under the same positive
denominator condition.
These rectangular and count estimates are developed in \cite{gtsharpmass}.

\subsection{Real brackets and full-time boundary signs}

\begin{proposition}[Lower mass from disjoint real brackets]
\label{gt-real-brackets}
Suppose the conditions of Proposition~\ref{gt-grouped-exterior} hold
throughout a closed time cell \(J\) for a real entire field \(H_t\).
Let \((b_j^-(t),b_j^+(t))\) be finitely many positive real intervals
outside the auxiliary disk, pairwise disjoint on each slice, with
continuous endpoints and strict opposite signs of \(H_t\) at their
endpoints for all \(t\in J\). Suppose \(c(t)\in[c_-,c_+]\),
\(D(t)\in[D_-,D_+]\), and put
\[
 l_j=\inf_J b_j^-,\qquad u_j=\sup_J b_j^+,\qquad
 U_j=\max\{|l_j-c_+|,|u_j-c_-|\}.
\]
The intervals supply at least one distinct real zero occurrence each
on every slice and
\begin{equation}
 -K(t)\ge\sum_j\frac{1}{U_j^2-D_-}\qquad(t\in J).
 \label{gt-real-bracket-mass}
\end{equation}
For an even field with \(c_->0,\ c_-^2>D_+\), if the reflected target
pair is exterior and disjoint from the positive brackets and their
reflections, one may add its lower mass \(1/[2(c_+^2-D_-)]\).
If the bracket reflections are also exterior, one may add
\(\sum_j[(u_j+c_+)^2-D_-]^{-1}\).
\end{proposition}

\begin{proof}
At each time the intermediate value theorem gives a real zero in
each bracket. Select one occurrence per interval on that slice.
Disjointness prevents duplicate selection. The remaining multiset
is still conjugation-invariant, even if further copies of a selected
real root remain. For a selected zero \(r_j\), the distance bound
comes from the extremes of \([l_j,u_j]-[c_-,c_+]\). Exterior
separation ensures \(0<(r_j-c)^2-D\le U_j^2-D_-\). Taking
reciprocals and omitting the favorable remainder proves
\eqref{gt-real-bracket-mass}. The reflected target pair contributes
\(1/[2(c^2-D)]\). Evenness gives one distinct negative reflection of
each selected real occurrence. For those exterior reflections the
positive denominator is at most \((u_j+c_+)^2-D_-\), because
\(r_j,c>0\). The disjointness assumptions allow each contribution
to enter once.
\end{proof}

The count is a lower count and the sign test allows multiple zeros.
Locations may be selected separately on each slice. In particular,
brackets with overlapping unions of locations over \(J\) still give
distinct occurrences if they are disjoint at the same time. An
independent positive lower bound \(D_-\) strengthens the weights and
already establishes reality of the target pair on that cell.

The boundary signs can be established on entire cells using the
physical heat equation. For an affine real boundary
\(b(t)=b_*+v(t-t_*)\) and any fixed positive factor \(\gamma\), set
\(f(t)=\gamma H_t(b(t))\). The operators \(\partial_t\) and
\(\partial_z\) commute, so \(\partial_t H_t=-\partial_z^2H_t\) gives
\begin{equation}
 f^{(m)}(t)=\gamma\sum_{j=0}^m\binom mj(-1)^jv^{m-j}
                     \partial_z^{m+j}H_t(b(t)).
 \label{gt-moving-boundary-derivative}
\end{equation}
If \(|\gamma H_t|\le M_*\) on every disk of radius \(R_*\) centered
at \(b(t)\), Cauchy's estimate bounds \(|f^{(m)}|\) by
\[
 C_m=M_*\sum_{j=0}^m\binom mj |v|^{m-j}
                     \frac{(m+j)!}{R_*^{m+j}} .
\]
For example, if an estimate in the local gauge
\(\exp(\pi b(t)/8)\) is available, the fixed gauge
\(\gamma=\exp(\pi b_*/8)\) multiplies that disk bound by at most
\(\exp(\pi |v|\max_{t\in J}|t-t_*|/8)\).

On a cell \(J=[t_*-\tau,t_*+\tau]\), interpolate sample values
with absolute error \(\varepsilon\) in this fixed gauge at the
\(n+1\) nodes \(t_*+\tau\cos(\pi j/n)\), \(0\le j\le n\).
The discrete cosine expansion of the interpolant of the sample
errors has ordinary Chebyshev coefficients bounded by
\(\varepsilon\) in degrees \(0,n\) and \(2\varepsilon\) in the
intermediate degrees, by the weighted discrete orthogonality with
half endpoint weights. Its uniform norm is at most \(2n\varepsilon\).
The node product in the scaled variable \(x=(t-t_*)/\tau\) is
\(2^{1-n}(x^2-1)U_{n-1}(x)\), with \(U_{n-1}\) the Chebyshev
polynomial of the second kind. Its norm is at most \(2^{1-n}\),
as is seen at \(x=\cos\theta\) from the product of
\(\sin\theta\) and \(\sin(n\theta)\). The ordinary interpolation
remainder for the true samples therefore bounds the error of the
computed interpolant \(p(x)=\sum_{j=0}^n a_jT_j(x)\) by
\[
 E=2n\varepsilon+
     \frac{C_{n+1}}{(n+1)!}\,2^{1-n}\tau^{n+1}.
\]
Here the polynomial coefficients represent exact interpolation of
the approximate samples. Any extra coefficient approximation errors
can be added in absolute value using \(|T_j|\le1\).
The estimate \(|T'_j|\le j^2\), from
\(T'_j=jU_{j-1}\), gives a sufficient sign test on each closed
subcell of radius \(\eta\) in \(x\) with midpoint \(x_0\):
\[
 \sigma p(x_0)>E+\eta\sum_{j=1}^n j^2|a_j|,\qquad \sigma\in\{-1,1\}.
\]
This follows by the mean value theorem and proves \(\sigma f>0\)
throughout the subcell. Affine ordering and exclusion inequalities
hold throughout a cell if they hold at both endpoints.
With independently enclosed target coordinates, these bounds yield
a uniform exterior lower-mass test \cite{gtbrackets}. The fixed
canonical application is recorded in
Appendix~\ref{sec:gt-source-certification}.

\subsection{An exact even quartic with negative drift and a collision}

Take any \(a>10\) and \(0\le t\le 1/10\). Let \(s=t-1/20\) and
\[
 F(t,z)=\e^{-s\partial_z^2}(z^2-a^2)^2
       =z^4-(2a^2+12s)z^2+a^4+4a^2s+12s^2.
\]
The terminating expansion satisfies \(\partial_tF=-\partial_z^2F\).
The even polynomial has order zero, and at \(s=0\) it has double
zeros at \(\pm a\), with \(F_{zz}(t,a)=8a^2\). The two squared
roots are \(a^2+6s\pm S\) with \(S^2=8s(a^2+3s)\). For
\(|s|\le1/20\) and \(a>10\),
\[
 |S|^2\le\tfrac25(a^2+\tfrac3{20})<\tfrac49 a^2,\qquad
 |6s\pm S|<\tfrac7{10}a .
\]
Both squared roots have positive real part. Their square roots
\(z_1,z_2\) in the right half-plane satisfy \(|z_i+a|\ge a\), so
\(|z_i-a|=|6s\pm S|/|z_i+a|<7/10\). The remaining roots are
their negatives. The disks of radii \(9/10\) and \(8/5\) about \(a\)
therefore count the same two roots, and the strip holds since
\(7/10<\sqrt{1-2t}\) throughout the slab.

Writing \(w_i=z_i-a\) and averaging the identities
\(2aw_i=6s\pm S-w_i^2\) gives
\[
 |c-a|\le \frac{79}{200a}<\frac7{100}.
\]
Also \(D=S^2/(z_1+z_2)^2\) and
\(|z_1+z_2|\ge2a-7/5\ge9a/5\), hence
\(|D|<100/729<1/5\). The sum \(z_1+z_2\) is positive real by
conjugation, so \(D\) has the sign of \(s\) away from zero.
The only exterior group is the reflected pair. Thus
\(K=-1/(2(c^2-D))<0\) on the slab. With \(v=1/10\) and
\(W(t)=\exp(4\int_t^v K)\), the exact gap equation yields
\[
 D(v)-2\int_0^v W(t)\,\dd t=W(0)D(0)<0 .
\]
The margin is zero if the cost starts at the double. This control
shares the order, strip, geometry and favorable sign, and fails the
strict incoming-cost condition. Its polynomial is unbounded on the
real axis, whereas a Fourier transform of a finite nonnegative
measure there is bounded by the source mass.

\subsection{Scope of the strip input}

The strip and the order hypotheses have separate roles. For instance,
the endpoint-defective positive finite square-source changes and the
bounded positive finite-prime orbit changes have order at most one
by the same double-exponential bound, but their zeros escape every
fixed horizontal strip. To see the distinction, let \(f_t\) denote
their heat-weighted half-line source. Suppose \(f_t^{(2j-1)}(0)\ne0\)
is the first nonzero odd endpoint derivative. This first term
survives the even heat multiplier. Integrating the cosine transform by
parts \(2j+2\) times yields
\[
 \int_0^\infty f_t(u)\cos(zu)\,\dd u
 =\frac{(-1)^j f_t^{(2j-1)}(0)}{z^{2j}}
    +O_Y(|z|^{-2j-2}) \qquad (|\Im z|\le Y).
\]
The differentiated half-line sources are integrable against
\(\e^{Yu}\), so the last integration remainder has this order.
Thus each fixed strip contains only finitely many zeros. On the
positive imaginary axis the first nonzero mode dominates the positive
tail of these sources. On
\([(\log y)/4,(\log y)/4+1]\), its density is bounded below by
\(\exp(-O_t(y))\), and \(\cosh(yu)\ge\e^{yu}/2\).
This gives logarithmic growth at least
\((y/4)\log y-O_t(y)\), ruling out exponential type. An entire
function of order at most one with finitely many zeros is a polynomial
times a linear exponential by Hadamard factorization. There must
therefore be infinitely many zeros, even though each strip contains
finitely many. This argument depends on the indicated endpoint
defect and tail assumptions. A changed source meeting the strip,
count and order hypotheses still obeys the grouped estimate.


The individual premise screen is recorded in \cite[\S6]{gtjoint}.
Euler-factor translates can share order one and the strip when their
additional initial zeros lie within its width; the required enlarged
disk and useful lower mass are then still untested. Reflected
super-Gaussian controls fail order at most one. The damped smooth-shape
and positive-primitive controls share the order bound but have no
established common strip or local lower-mass payment here. Finite-profile
threshold controls do not supply this earlier slab merely by being
all-real at their threshold. The normalized endpoint-merger control
retains its RH hypothesis; its remote added doubles do not identify
this tube. Changed observations need not satisfy the physical heat
equation, and fractional profile masses do not supply an integer
multiplicity-counted divisor. Each conclusion concerns the named
premise and control class, not arbitrary changes of the theta source.
